\chapter{Základné pojmy}

Organizácie potrebujú na vykonávanie činností a dosiahnutí svojich úloh spracovávať informáciu. V neustále zväčšujúcej
sa miere je zaznamenávaná a uchovávaná v digitálnej forme, pričom je spracovaná pomocou digitálnych a informačných komunikačných
technológií (\acrshort{IKT}, \acrshort{dIKT}). 
Informatizácia spoločnosti je jednou z priorít Európskej únie\footnote{Digitálne desaťročie Európy: digitálne ciele na 
rok 2030}(\acrshort{EU}) a aj Slovenskej republiky\footnote{Stratégia digitálnej transformácie Slovenska 2030}. Informatizácia
pomáha efektívne a rýchlejšie spracovávať informácie, zvyšuje dostupnosť (sprístupnením cez internet za použitia mobilu
alebo osobného počítača) a umožňuje poskytovať občanom tradičné a nové služby.

Už v roku 2010, krátko po veľkej hospodárskej kríze, \acrshort{EU} zdôraznila význam \acrshort{dIKT}\footnote{V dokumente
Digitálna agenda pre európu: https://eur-lex.europa.eu/legal-content/sk/ALL/?uri=CELEX:52010DC0245} pre hospodársky rozvoj
členských štátov \acrshort{EU}. Bez spoľahlivého fungovania \acrshort{dIKT} nedokáže spoločnosť, jej inštitúcie ani jej
organizácie vykonávať svoje činnosti v potrebnom rozsahu a kvalite. V roku 2015 publikovala \acrshort{EU} smernicu 
NIS(z angl. \textit{Network and information systems}), kde boli \acrshort{dIKT} kategorizované ako prvok kritickej
infraštruktúry spoločnosti. Sformovala sa nová oblasť činnosti: informačná, resp. \acrfull{KIB}, ktorá sa zoberá ochranou 
zvlášť \acrshort{IKT} a zvlášť \acrshort{dIKT}.

Oblasť KIB je nová a dynamicky sa rozvíjajúca oblasť, v ktorej sa terminológia ešte neustálila\footnote{\acrshort{ENISA}
publikovala štúdiu, v ktorej porovnávala rozlične výklady pojmu kybernetická bezpečnosť naprieč šiestimi rôznymi inštitúciami.}.
V tejto kapitole definujeme základné pojmy \acrshort{KIB} a úlohy, ktoré rieši. Budeme vychádzať z úvodnej kapitoly dokumentu
\cite{KomentarZoKB}. Ak je čitateľ znalý v problematike \acrshort{KIB} môže pokojne prejsť k ďalšej kapitole. Ak by mal čitateľ
záujem o podrobnejšie informácie, dávame do pozornosti krátky úvod do problematiky \acrshort{KIB}\cite{KUDIKB-MVS}.

\begin{itemize}

\item \textbf{aktívum} je čokoľvek, čo má pre organizáciu hodnotu, môžu byť cieľom útoku a má zmysel (resp. je potrebné) chrániť
pred potenciálnymi hrozbami. Aktíva sú hmotné (zariadenia, infraštruktúra, personál) a nehmotné (peniaze, informácie, vedomosti).

\item \textbf{hrozba} je objektívne existujúca potenciálna možnosť priamo alebo nepriamo narušiť (alebo negatívne ovplyvniť) 
systém, spracovávané informácie, alebo iné aktíva organizácie.

\item \textbf{\acrfull{KIB}} skúma hrozby voči aktívam a hľadá opatrenia, ktoré majú
eliminovať riziká vyplývajúce z hrozieb voči aktívam. Táto disciplína presadzuje metódy dlhodobej, nákladovo a funkčne efektívnej
ochrany aktív organizácie.

\item \textbf{systém riadenia informačnej bezpečnosti}, angl. \acrfull{ISMS}, je súbor nástrojov, metód a postupov pre zaistenie
systematického riešenia informačnej bezpečnosti v organizácií. Cieľom tohto systému je zaručiť dôvernosť, integritu a dostupnosť
informácie, procesov a systémov informačných technológií. Podrobnejšie sa ním budeme zaoberať v kapitole 4 o zavádzaní ISMS.

\end{itemize}
